% © 2026 Ing. David Jaroš — CC BY-NC-ND 4.0
%
% This work is licensed under a Creative Commons Attribution-NonCommercial-NoDerivatives
% 4.0 International License (CC BY-NC-ND 4.0).
%
% License History: Earlier drafts (up to v0.3) were released under CC BY 4.0.
% From v0.4 onward, all material is released under CC BY-NC-ND 4.0 to protect
% the integrity of the theoretical work during ongoing academic development.
%
% See LICENSE.md for full license text.

\ifdefined\INCLUDEMODE
\else
  \documentclass[12pt,a4paper]{article}
  \usepackage[a4paper, margin=2.5cm]{geometry}
  \usepackage{amsmath, amssymb, amsthm}
  \usepackage{mathtools}
  \usepackage{hyperref}
  \usepackage{xcolor}
  \usepackage{mdframed}
  \usepackage{booktabs}

  \newtheorem{theorem}{Theorem}[section]
  \newtheorem{lemma}[theorem]{Lemma}
  \newtheorem{proposition}[theorem]{Proposition}
  \theoremstyle{definition}
  \newtheorem{definition}[theorem]{Definition}
  \theoremstyle{remark}
  \newtheorem{remark}[theorem]{Remark}

  \newmdenv[backgroundcolor=yellow!20, linecolor=orange, linewidth=1.5pt]{gapbox}
  \newmdenv[backgroundcolor=green!15, linecolor=green!60!black, linewidth=1.5pt]{resultbox}
  \newmdenv[backgroundcolor=blue!8, linecolor=blue!50, linewidth=1.5pt]{insightbox}

  \title{\textbf{Step 7: Non-Commutativity in the Full Biquaternionic FPE}\\
         \large Gap G3: Does $[\partial_T, \nabla_Q^2]$ Vanish in $\mathbb{C}\otimes\mathbb{H}$?}
  \author{David Jaro\v{s}}
  \date{2026-03-06}

  \begin{document}
  \maketitle

  \begin{abstract}
  The FPE equivalence (Step~4) was proved in the scalar sector.
  Gap~G3 asks whether non-commutativity in the full biquaternionic FPE
  generates new terms or cancels by the algebraic structure of $\mathbb{C}\otimes\mathbb{H}$.
  We show that since $\mathbb{C}\otimes\mathbb{H}$ is associative
  (proved), the FPE remains form-invariant under the natural ordering
  prescription, but a residual operator-ordering ambiguity remains for
  the drift term $A(Q)\Theta$ when $A$ and $\Theta$ are non-commuting
  biquaternions.  Gap G3 is reduced but not closed.
  \end{abstract}
\fi

%──────────────────────────────────────────────────────────────────────────────
\section{Context: The Biquaternionic FPE}
\label{sec:context}
%──────────────────────────────────────────────────────────────────────────────

The biquaternionic Fokker-Planck equation (Step~4) is:
\begin{equation}
  \partial_T\Theta \;=\; -\nabla_Q\cdot\bigl[A(Q)\Theta\bigr]
                         + D\,\nabla_Q^2\Theta,
  \label{eq:FPE}
\end{equation}
where $T \in \mathbb{C}$ (complex time), $Q \in \mathbb{H}$ (quaternionic
position), $A(Q) \in \mathbb{C}\otimes\mathbb{H}$ (drift), and
$D \in \mathbb{R}^+$ (diffusion).

In the scalar sector (where $A$ is real and $\Theta$ commutes with $A$),
the FPE is proved to be the Euler-Lagrange equation of $S[\Theta]$
(Step~4, Theorem~1).  The question is whether this holds in the full
biquaternionic case.

%──────────────────────────────────────────────────────────────────────────────
\section{Non-Commutativity Analysis}
\label{sec:noncomm}
%──────────────────────────────────────────────────────────────────────────────

\begin{definition}[Left and Right FPE conventions]
In $\mathbb{C}\otimes\mathbb{H}$, the product $A\Theta \neq \Theta A$ in general.
Two natural conventions:
\begin{align}
  \text{Left FPE:}   &\quad \partial_T\Theta = -\nabla_Q\cdot[A\cdot\Theta] + D\nabla_Q^2\Theta, \\
  \text{Right FPE:}  &\quad \partial_T\Theta = -\nabla_Q\cdot[\Theta\cdot A] + D\nabla_Q^2\Theta.
\end{align}
\end{definition}

\begin{lemma}[Commutator $[\partial_T, \nabla_Q^2]$ in $\mathbb{C}\otimes\mathbb{H}$]
\label{lem:commutator}
Since $\partial_T$ acts on the complex-time component and $\nabla_Q^2$
acts on the quaternionic-position component, and since
$\mathbb{C}\otimes\mathbb{H}$ is an algebra over $\mathbb{R}$:
\begin{equation}
  [\partial_T, \nabla_Q^2]\,\Theta \;=\; 0
  \quad\text{(form-invariant)}.
\end{equation}
\end{lemma}

\begin{proof}
$\partial_T = \partial_{t} + i\partial_\psi$ acts on the
$\mathbb{C}$-factor of $\mathbb{C}\otimes\mathbb{H}$.
$\nabla_Q^2 = \partial_{Q_0}^2 + \partial_{Q_1}^2 + \partial_{Q_2}^2 + \partial_{Q_3}^2$
acts on the $\mathbb{H}$-factor.
Since $T$ and $Q$ are independent coordinates:
\begin{equation}
  \partial_T\bigl(\nabla_Q^2\Theta\bigr)
  \;=\; \nabla_Q^2\bigl(\partial_T\Theta\bigr),
\end{equation}
i.e., $[\partial_T, \nabla_Q^2] = 0$ on smooth $\Theta$.
\end{proof}

\begin{insightbox}
\textbf{Conclusion from Lemma~\ref{lem:commutator}:}
The kinetic (diffusion) part of the FPE is \emph{form-invariant} under
the full biquaternionic structure.  The non-commutativity does not
generate new terms in the $D\nabla_Q^2\Theta$ part.
\end{insightbox}

%──────────────────────────────────────────────────────────────────────────────
\section{Residual Ambiguity: Drift Term Ordering}
\label{sec:drift}
%──────────────────────────────────────────────────────────────────────────────

The issue in Gap~G3 is the drift term $-\nabla_Q\cdot[A(Q)\Theta]$ when
$A, \Theta \in \mathbb{C}\otimes\mathbb{H}$ and $[A, \Theta] \neq 0$.

The Euler-Lagrange equation derived from $S[\Theta]$ (with drift
$A = -\nabla_Q\mathbb{H}$) gives:
\begin{equation}
  \frac{\delta S[\Theta]}{\delta\bar{\Theta}}
  \;=\; \partial_T\Theta + \nabla_Q\cdot[(\nabla_Q\mathbb{H})\Theta] - D\nabla_Q^2\Theta
  \;=\; 0.
\end{equation}
In the biquaternionic case, the product $(\nabla_Q\mathbb{H})\Theta$ uses
\emph{left} multiplication.  This selects the Left FPE convention.

\begin{proposition}[G3 partial resolution: Left FPE from $S[\Theta]$]
\label{prop:G3}
The Euler-Lagrange equations of $S[\Theta]$ give the \textbf{left} FPE
convention.  The right FPE is not an Euler-Lagrange equation of any action
of the same form.  Therefore the ordering is fixed: left multiplication for the drift.
\end{proposition}

\begin{proof}
The action $S[\Theta] = \int \langle D_T\Theta, D_T\Theta\rangle\,dT\,dQ$
with $D_T = \partial_T + A(Q)$ uses the \emph{left} covariant derivative
$A(Q)\Theta$.  Varying $\delta S/\delta\bar{\Theta} = 0$ gives the
left FPE.  A right-convention action would require $\Theta\cdot A(Q)$
in the covariant derivative, which breaks the gauge structure of $S[\Theta]$
(since gauge fields act from the left in the UBT covariant derivative).
\end{proof}

%──────────────────────────────────────────────────────────────────────────────
\section{Remaining Issue}
\label{sec:remaining}
%──────────────────────────────────────────────────────────────────────────────

\begin{gapbox}
\textbf{Gap G3: Partially resolved.}

\textbf{Resolved:} $[\partial_T, \nabla_Q^2] = 0$ (Lemma~\ref{lem:commutator}).
The diffusion term is form-invariant.

\textbf{Resolved:} Ordering is fixed as left FPE (Prop.~\ref{prop:G3}).

\textbf{Remaining:}
In the Hamiltonian sector $A = -\nabla_Q\mathbb{H}$, the biquaternionic
gradient $\nabla_Q\mathbb{H}$ may itself have ordering ambiguity when
$\mathbb{H}$ (the Hamiltonian, not the quaternions) is a biquaternionic-valued
function.  Specifically: $\nabla_Q(\mathbb{H}\Theta)$ vs $(\nabla_Q\mathbb{H})\Theta
+ \mathbb{H}(\nabla_Q\Theta)$ --- the product rule must be checked for
non-commuting $\mathbb{H}$ and $\Theta$.

By the Leibniz rule in $\mathbb{C}\otimes\mathbb{H}$ (which is a differentiable
algebra), this product rule holds for left multiplication.
Therefore the remaining ordering issue reduces to:
is $\mathbb{H}$ always left-multiplying $\Theta$?
In UBT: yes, because $\mathbb{H}$ is the Hamiltonian (a scalar-valued
energy) and $\Theta$ is the biquaternionic field.

\textbf{Conclusion:} Gap G3 is effectively closed for the physically
relevant case where $\mathbb{H}$ is a scalar (or diagonal in $\mathbb{H}$).
Full closure requires checking the vector potential case ($A_\mu \in
\mathfrak{su}(2) \subset \mathbb{H}$) separately.
\end{gapbox}

\ifdefined\INCLUDEMODE
\else
  \end{document}
\fi
